% Standalone paragraph. Paste into your own thesis, or \input it.
% Not wired into any chapter's numbering -- drop it wherever the "the tree has to be extracted
% before it can be measured" argument is needed. Companion to whole_tree_gap.tex,
% reproducible_features.tex and healthy_cohort_2.tex: those argue the field measures the wrong
% place in the tree, cannot agree on how to measure it, and measures the wrong people; this one
% argues that the tree handed to all three is not reliably the tree that is there.

% HEADLINE: one of four sibling subsections for the State of the Art chapter, sitting under
% \section{The gap} (state_of_the_art/08_the_gap.tex): how the tree is extracted
% (segmentation_gap), where it is measured (whole_tree_gap), how it is measured
% (reproducible_features), in whom it is measured (healthy_cohort_2). Promote to \section if they
% are to stand as sections of their own.

% VERIFICATION STATUS of the citations used here, so a later pass knows what each can bear:
%   bransby2026imagecasx        -- read in full. Every figure below was read out of Table 2 of the
%                                  PDF at the repo root, and the vessel-break clause is quoted from
%                                  its Fig. 6 discussion.
%   medranogracia2016atlas      -- [BIBLIOGRAPHY VERIFIED, READ FIRST-HAND] for exactly the 3%
%                                  mesh-failure claim used here, and for nothing beyond it.
%   dong2023casnet, zhang2024adehtl, isensee2021nnunet, shit2021cldice, qiu2025topology,
%   nannini2024characterization -- [BIBLIOGRAPHY VERIFIED] only, i.e. abstract level. No sentence
%                                  here claims more than an abstract supports; re-check each
%                                  against the full text before submission.

% UNRESOLVED, ask the dataset authors: the paper gives the inter-observer Betti error as 0.2 in its
% Results text and in Table 1 ("All segments"), but 0.4 in Table 2. This file quotes Table 2,
% because that is the table the automated methods are scored in. If 0.2 is the right value, the
% ratio in the opening paragraph becomes ten-fold rather than five-fold.

% QUESTIONS FOR SUPERVISOR (working notes, not thesis text -- delete before pasting):
%  1. Reproduce or train? CAS-Net, 3D-FFR-UNet and Swin UNETR run today from the delivered weights,
%     as staged run directories that need no config change. Does reproducing the published
%     benchmark satisfy objective 5, or does the thesis have to train a segmentation model itself?
%  2. Selected on which metric? CAS-Net leads on DSC and Betti error; ADE-HTL ties it on every
%     centerline metric and would need TotalSegmentator heart-chamber masks precomputed before it
%     can run here at all. If the downstream unit of analysis is the tree, should the segmenter be
%     chosen on Betti error and clDice rather than DSC -- and is that worth the extra dependency?
%  3. Inside the network or after it? This framework has no clDice or soft-skeleton loss; clDice
%     exists here only as an evaluation metric, and ADE-HTL is the only method in it that makes
%     topology a training signal. Given that the benchmark's own nnU-Net+clDice row got worse on
%     Betti error, do we accept the delivered masks, fine-tune with a topology-aware loss
%     (ThesisPlan.md week 5, currently conditional), or add a post-hoc reconnection step that
%     leaves the network untouched?
%  4. Validating on CGPS, where no coronary ground truth exists yet: is the ImageCAS-trained model
%     applied as is, is the hand-labelled subset the only accuracy check, and can repeat-scan
%     agreement stand in for accuracy where labels will never exist?

\subsection{How the tree is extracted}
\label{sec:gap-obtaining}

Automated segmentation now matches expert annotators on how much of the artery it finds, and does
not match them on whether what it finds is connected.
The benchmark published with the dataset used here scored eight methods against two analysts on the
same 160 coronary computed tomography angiography (CCTA) scans.
Its best method, CAS-Net, reached a Dice similarity coefficient (DSC, the share of voxels the
prediction and the reference annotation have in common) of 91.2 against 92.8 for the analysts, and a
Betti error --- the number of connected components and loops by which the predicted tree differs
from the true one --- of 1.9 against their 0.4 \cite{bransby2026imagecasx, dong2023casnet}.
The gap in overlap is 1.6 points; the gap in topology is a factor of nearly five.
All of these figures come from that one evaluation.

Overlap cannot see that difference.
The alternative, clDice, scores agreement along a differentiable skeleton of the vessel rather than
over its volume,
and was introduced with a proof that it preserves topology up to homotopy equivalence
\cite{shit2021cldice}: two masks can score the same DSC and carry different trees.

Optimizing that metric does not fix what it measures.
In the same benchmark, nnU-Net trained with an added clDice term gained 0.2 DSC over plain nnU-Net
and returned the worst Betti error in the table, 8.0 \cite{isensee2021nnunet}.
The two leading methods separate on overlap and tie on connectivity: CAS-Net beats ADE-HTL on DSC
and Betti error (both $p<0.001$) but is indistinguishable from it on clDice ($p=0.94$) and on both
centerline distance measures ($p=0.74$, $p=0.90$) \cite{zhang2024adehtl}.
The benchmark's own reading of its qualitative examples is that vessel breaks are present in all
model predictions despite high DSC and clDice \cite{bransby2026imagecasx}.

Repairing the tree after the network is where recent work puts its effort.
Qiu et al.\ combine a centerline-enhanced loss, a reconnection step that matches broken ends by
distance and direction, and implicit reconstruction of vessels missing altogether, reaching a DSC of
88.5 and 85.1 on two other public datasets, ASOCA and PDSCA \cite{qiu2025topology}.
Those are not the cohort above, so the numbers cannot be set beside it; what the work establishes is
that a repair stage is still thought necessary after a well-trained network.

A break costs little overlap and can cost a whole feature: losing a distal branch moves DSC by a
fraction of a point, removes a branch from a branch count, and in the posterior vessels can change
which artery is called dominant.
The resolution limit behind such breaks also defeats careful manual work --- in an atlas of normal
anatomy built from 300 scans selected for zero calcium and no stenosis, 3\% of cases could not be
meshed at all \cite{medranogracia2016atlas}.

Segmentation is therefore inherited here rather than rebuilt: trained weights ship with the
benchmark, so tree extraction is reproduced \cite{bransby2026imagecasx}, and the closest published
pipeline --- segmentation, centerline extraction and geometric description run end to end on 281
single-center scans --- stops at description \cite{nannini2024characterization}.
What remains open is which segmentation the tree is built from, whether topological correctness is
trained into the network or repaired after it, and how either is checked in a cohort that carries no
reference annotation at all.
